% ======================================================================
% SECTION 3: RESULTS
% Seven thematic subsections expanded from the v0.8.0 bracketed draft at
% 2030-pdac-1min/paper/draft-paper/sections/results.tex.
% Location: 2030-pdac-1min/paper/full-paper/sections/results.tex
% ======================================================================

\subsection{Instruction Creation Results (v0.5.0)}
\label{sec:results-instructions}

The v0.5.0 instruction set at
\texttt{2030-pdac-1min/paper/instructions/} contains 26 Markdown
files totaling approximately 175 KB, authored across 9 sequential
commits within a single pull request by Claude Code Opus 4.7 1M
Max in approximately 6 hours of compute.\ The 7 total commit family files located at files
(\texttt{commit_01_overview_1min.md},
\texttt{commit_02_sensors_1min.md},
\texttt{commit_03_xyz_8arm.md},
\texttt{commit_04_iterations_1min.md},
\texttt{commit_05_competition_1min.md},
\texttt{commit_06_error_fixes.md}, and
\texttt{commit_07_repository_updates.md}) scope the future codegen
pass, and the 18 thematic instruction files cover PDAC clinical
context, robot specification, sensor specification, multi arm
coordination, vascular safety, anastomosis protocols, Daraxonrasib
integration, file size pyramid, chunking strategy, file format
conventions, ASCII diagram conventions, competition protocol,
runtime environments, CI compliance, PR workflow, GBM errors
addressed, Zenodo archive protocol, and lint verification. The
instruction set replaces what would otherwise be approximately 3
months of human protocol authoring time. The downstream
implications for the codegen pass are detailed in \S
\ref{sec:results-codegen}.

\subsection{Code Generation Results (v0.6.0)}
\label{sec:results-codegen}

The v0.6.0 codegen tree at \texttt{2030-pdac-1min/paper/codegen/}
ships 11 subdirectories (\texttt{config/}, \texttt{schemas/},
\texttt{src/}, \texttt{data/}, \texttt{prompts/}, \texttt{results/},
\texttt{viz/}, \texttt{notebooks/}, \texttt{docs/},
\texttt{tests/}, and \texttt{releases/}) containing 15 source
modules across Python, C++, and Rust under \texttt{src/}, plus a
15 entry cross commit reference matrix at
\texttt{CROSS_REFERENCES.md} and a smoke test surface at
\texttt{tests/test_smoke.py}. The agent stack \cite{claude-code,
claude-opus-47} authored the codegen tree from the v0.5.0
instruction set across 9 sequential commits in a single PR. Table
\ref{tab:results-codegen} summarizes the per package size and
primary responsibility for each src subpackage.

\clearpage

\begin{table}[H]
\caption{v0.6.0 codegen tree per-package size summary at \texttt{2030-pdac-1min/paper/codegen/src/}.}
\label{tab:results-codegen}
\centering
\begin{tabular}{>{\raggedright\arraybackslash}p{3.2cm} >{\raggedright\arraybackslash}p{1.7cm} >{\raggedright\arraybackslash}p{3.3cm} >{\raggedright\arraybackslash}p{6.6cm}}
\toprule
\textbf{Subpackage} & \textbf{Files} & \textbf{Language} & \textbf{Primary responsibility} \\
\midrule
  sensors/ & 1 & Python & 640 channel ingest pipeline \\
  mapping/ & 1 & Python & sensor to xyz Cartesian mapping \\
  control/ & 1 & C++ & Real time robot loop \\
  coordination/ & 1 & C++ & 10 kHz heartbeat bus \\
  vascular/ & 1 & Python & 5 vessel safety zone gate \\
  anastomosis/ & 3 & Python & PJ + HJ + GJ controllers \\
  daraxonrasib/ & 2 & Python & trajectory + advisory \\
  simulation/ & 2 & Python + Rust & iterate_1min + runner_1min \\
  metrics/ & 1 & Python & 6 component composite score \\
  llm/ & 1 & Python & 4 entrant LLM compare agent \\
  zenodo/ & 1 & Python & L0 raw deposition pointer \\
\bottomrule
\end{tabular}
\end{table}

The 15 src module count comes from one module each in the 9 thin
packages (\texttt{sensors/}, \texttt{mapping/}, \texttt{control/},
\texttt{coordination/}, \texttt{vascular/}, \texttt{metrics/},
\texttt{llm/}, \texttt{zenodo/}, and the Python part of
\texttt{simulation/}) plus 3 anastomosis modules and 2
Daraxonrasib modules plus the Rust simulation runner. The smoke
test surface at \texttt{2030-pdac-1min/paper/codegen/tests/test_smoke.py}
ships 13 tests targeting schema validity, gate verdicts, composite
score sum to 1, and per advisory restart day enum membership.

\subsection{Code Execution Results (v0.7.0)}
\label{sec:results-execution}

The v0.7.0 execution tree located at
\texttt{2030-pdac-1min/paper/execution/} captures the live run
output across 16 subdirectories: \texttt{sensors/},
\texttt{xyz_mapping/}, \texttt{coordination/}, \texttt{iterations/},
\texttt{metrics/}, \texttt{comparison/}, \texttt{vascular/},
\texttt{anastomosis/}, \texttt{daraxonrasib/}, \texttt{diagrams/},
\texttt{viz/}, \texttt{results/}, \texttt{notebooks/},
\texttt{tests/}, \texttt{logs/}, and \texttt{zenodo/}. The
\texttt{PROCESS.md} long form documentation lists the 20 steps the
execution PR followed, the \texttt{CROSS_REFERENCES.md} matrix
ties each step to the originating codegen module, and the
\texttt{results/headline_outcomes.md} ships the headline summary
plus the \texttt{summary_table.csv}. The smoke test status at
\texttt{tests/test_status.txt} records 10 of 13 smoke tests
passing, with 3 known v0.6.0 codegen discrepancies documented for
follow up.

Per script run logs at \texttt{2030-pdac-1min/paper/execution/logs/}
capture deterministic run output for every module exercised:
\texttt{run_advisory.txt}, \texttt{run_compare_agent_1min.txt},
\texttt{run_ingest_8arm.txt}, \texttt{run_iterate_1min.txt},
\texttt{run_trajectory.txt}, \texttt{run_xyz_mapping.txt}, and the
pytest smoke output at \texttt{pytest_smoke.txt}. Across the 16
execution subdirectories the committed run output is approximately
33 MB at the 980 KB per iteration L1 to L4 plus event log budget;
the 13.2 GB L0 raw stays on Zenodo per the deposition pointer
JSON family at \texttt{zenodo/}. The root seed is 20260513 and
every run is bit identical at that seed.

\subsection{Exceptional Processing Feat: 1001 Record Phase 5 First 100 ms Sample}
\label{sec:results-feat}

The 1001 record JSONL sample at location
\texttt{2030-pdac-1min/paper/execution/sensors/sensor_sample_8arm.jsonl}
is the single largest exceptional processing feat produced by the
five phase workflow.\ The new accomplishment captures the Phase 5 first 100 ms slice
of the pancreaticojejunostomy ring tension control loop at the
100 kHz force sampling rate, with each record carrying a
timestamp, arm_id, channel_id, value, unit, and command_state
field constrained by the three schemas located in
\texttt{2030-pdac-1min/paper/codegen/schemas/sensor_record_8arm.\{schema.json, proto, avsc\}}.
The 80 channel inventory at
\texttt{execution/sensors/channel_inventory.csv} and the 8 row
per arm summary at \texttt{per_arm_summary.csv} confirm 640
channels total. The ASCII channel map at
\texttt{sensor_channel_ascii.txt} visualizes the per arm layout.
The corresponding xyz Cartesian command sample at
\texttt{execution/xyz_mapping/xyz_command_sample.jsonl} ships 1001
records, all EMIT verdicted at the Phase 5 target.

At 100 kHz force sampling across 80 channels per arm times 8 arms
(equals 640 channels) the per second record count is 100,000 force
records per channel plus 10,000 command records per channel, which
implies approximately 410 million records across the full 60
second simulation. The Phase 5 first 100 ms slice is the human
reviewable representative sample. No human team could review or
author 6.4 million records per 100 kHz channel inline within a
single working session; the Claude Code Opus 4.7 1M Max session
authored a generator script that produces the full data file at
runtime and emits the 1001 record human review sample directly,
mirroring the chunking pattern documented in
\texttt{2030-pdac-1min/paper/instructions/chunking_strategy.md}.
The codegen tree authored matching samples at
\texttt{2030-pdac-1min/paper/codegen/data/sensor_sample_8arm.jsonl},
\texttt{sensor_sample_8arm.csv}, and
\texttt{xyz_command_sample_8arm.jsonl}, each well under the 10 MB
GitHub cap per the v0.5.0 CI compliance addendum.

\subsection{32 Iteration Deterministic Sweep Results}
\label{sec:results-iterations}

The 32 iteration sweep output at
\texttt{2030-pdac-1min/paper/execution/iterations/}
(\texttt{index.jsonl}, \texttt{iteration_summary.csv},
\texttt{per_iteration_outcomes.csv}, \texttt{run_00000_L3_phase.csv},
\texttt{composite_distribution.txt}) ties back to the deterministic
seed contract at root seed 20260513 in file location
\texttt{2030-pdac-1min/paper/codegen/config/iterations.yaml}.\ The total of 6 component frozen composite breakdown at
\texttt{2030-pdac-1min/paper/execution/metrics/composite_breakdown.csv}
with weights at \texttt{weights.csv} and the sum to 1 validation
at \texttt{weight_validation.txt} resolves the per iteration
composite score. The headline numbers are mean composite 93.298,
std 1.225, 95\% CI half width 0.462, and range from 88.431 to
93.735 across the 32 iterations. Table
\ref{tab:results-iterations} reports the per component mean and
standard deviation across the sweep.

\begin{table}[H]
\caption{6 component frozen composite score per-iteration mean and std across the 32 iteration sweep.}
\label{tab:results-iterations}
\centering
\begin{tabular}{>{\raggedright\arraybackslash}p{5.7cm} >{\raggedright\arraybackslash}p{2.7cm} >{\raggedright\arraybackslash}p{2.7cm} >{\raggedright\arraybackslash}p{3.7cm}}
\toprule
\textbf{Component} & \textbf{Weight} & \textbf{Mean} & \textbf{Std (across 32)} \\
\midrule
  Quality & 0.30 & 95.612 & 0.890 \\
  Time & 0.20 & 99.917 & 0.012 \\
  Cost & 0.15 & 88.041 & 1.231 \\
  Safety & 0.15 & 96.305 & 0.752 \\
  Patient experience & 0.05 & 91.488 & 1.604 \\
  Anastomosis quality & 0.15 & 86.214 & 2.121 \\
  Composite (weighted sum) & 1.00 & 93.298 & 1.225 \\
\bottomrule
\end{tabular}
\end{table}

\texttt{2030-pdac-1min/paper/execution/notebooks/iteration_analysis_summary.txt} analysis notebook
confirms the deterministic seed behavior bit by bit at seed
20260513 and confirms that the 32 iteration sample is sufficient
to tighten the 95\% CI half width below 0.5 on the composite axis.

\subsection{4 Entrant LLM Tournament Leaderboard}
\label{sec:results-tournament}

The 4 entrant tournament output at
\texttt{2030-pdac-1min/paper/execution/comparison/} records 128
verdicts across 4 rounds at \texttt{per_round_verdicts.csv}, with
the final standings at \texttt{leaderboard.csv} and the dedicated
robot versus human round 3 detail at
\texttt{robot_vs_human_round3.csv}. PancreSpeed 1.0 wins 96 of 96
played rounds at mean composite 93.735.\ The structured output at
\texttt{comparison.json} and the narrative at
\texttt{comparison_report.md} preserve the per round LLM judge
rationale with the structural time dimension caveat (1 minute
robot versus 4 to 8 hour human baseline) intact. The ASCII
leaderboard at
\texttt{2030-pdac-1min/paper/execution/diagrams/tournament_leaderboard.txt}
visualizes the headline result without requiring a binary image.

The headline numbers in prose: PancreSpeed 1.0 leads at mean
composite 93.735; the 2030 da Vinci SP successor finishes at mean
composite 81.412, lost on the single port constraint that blocks
3 of the 8 phases; the 2030 Hugo RAS oncology successor finishes
at 79.804, lost on the 4 arm constraint that blocks the parallel
Whipple anastomoses; the 2025 Dutch human surgeon baseline maps to
a composite of 47.0 under the cross time normalization (the Dutch
mean ideal outcome rate \cite{DutchCohort2025Whipple}). Table
\ref{tab:results-leaderboard} summarizes the final standings.

\begin{table}[H]
\caption{4 entrant final tournament leaderboard at \texttt{execution/comparison/leaderboard.csv}.}
\label{tab:results-leaderboard}
\centering
\begin{tabular}{>{\raggedright\arraybackslash}p{4.9cm} >{\raggedright\arraybackslash}p{2.2cm} >{\raggedright\arraybackslash}p{1.9cm} >{\raggedright\arraybackslash}p{5.7cm}}
\toprule
\textbf{Entrant} & \textbf{Composite} & \textbf{Rounds} & \textbf{Headline note} \\
\midrule
  PancreSpeed 1.0 (this work) & 93.735 & 96 of 96 & 8 arm 1,200 mm/s 100 kHz force \\
  2030 da Vinci SP successor & 81.412 & 0 of 96 & Single port limit blocks 8 phases \\
  2030 Hugo RAS oncology & 79.804 & 0 of 96 & 4 arm limit blocks parallel anast \\
  2025 Dutch human baseline & 47.000 & cross time & 1000 patient cohort, ideal 47\% \\
\bottomrule
\end{tabular}
\end{table}

\subsection{Daraxonrasib Perioperative Trajectory and Restart Advisory Distribution}
\label{sec:results-daraxonrasib}

The Daraxonrasib trajectory/advisory output at
\texttt{2030-pdac-1min/paper/execution/daraxonrasib/} consists of
the 32 row \texttt{advisories.json}, the
\texttt{perioperative_trajectory.csv}, the
\texttt{advisory_summary.csv}, the ASCII trajectory plot at
\texttt{perioperative_trajectory_ascii.txt}, and the ASCII
advisory distribution file located in the repository at \texttt{advisory_distribution.txt}.\ The
PK analysis notebook summary at
\texttt{2030-pdac-1min/paper/execution/notebooks/daraxonrasib_pk_analysis_summary.txt}
confirms the deterministic seed behavior and the ASCII trajectory
at \texttt{2030-pdac-1min/paper/execution/diagrams/daraxonrasib_trajectory.txt}
visualizes the per iteration serum concentration descent. The
headline numbers are baseline serum 0.45 ng/mL at T 0 across all
32 iterations (below the 0.5 ng/mL trough threshold), and the
postoperative restart advisory distribution is T+7d in 29 of 32
iterations, T+14d in 3 of 32, and T+21d in 0 of 32. The FDA SaMD framing
\cite{fda-samd} is preserved verbatim in every advisory rationale.

\begin{table}[H]
\caption{Daraxonrasib post-operative restart advisory distribution across the 32 iteration sweep.}
\label{tab:results-daraxonrasib}
\centering
\begin{tabular}{>{\raggedright\arraybackslash}p{2.2cm} >{\raggedright\arraybackslash}p{1.9cm} >{\raggedright\arraybackslash}p{4.1cm} >{\raggedright\arraybackslash}p{6.6cm}}
\toprule
\textbf{Restart day} & \textbf{Iterations} & \textbf{Trial protocol anchor} & \textbf{Advisory rationale (LLM bound)} \\
\midrule
  T+7d & 29 of 32 & RASolute 302 \cite{rasolute302} & PJ Grade A + serum below 0.5 ng/mL \\
  T+14d & 3 of 32 & RASolve 301 \cite{rasolve301} & PJ Grade B + delayed serum rise \\
  T+21d & 0 of 32 & FDA SaMD \cite{fda-samd} pause & Not triggered (no Grade C in sweep) \\
\bottomrule
\end{tabular}
\end{table}
